\section{A Fourier-increment class beyond tanh}
\label{app:fourier-activations}

The scalar argument extends beyond tanh. A mixed-activation normal form
lets us retain a short prefix of the original activation and use tanh
only in the replacement suffix. This separates the spectral condition
needed for the estimator from the geometric conditions needed for the
normal form.

\begin{theorem}[A bounded Fourier-increment activation class]
\label{thm:fourier-class}
Fix a real activation $\phi$ satisfying the following conditions.
\begin{enumerate}[label=(\roman*)]
\item $\phi(0)=0$, $\phi$ is one-Lipschitz, and $|\phi|\le A$ for
some fixed finite $A$.
\item For fixed $c,r_0,\kappa>0$,
$|\phi(u)-u|\le c|u|^{1+\kappa}$ when $|u|\le r_0$.
On an interval about zero it has a continuous inverse with
$|\phi^{-1}(v)|\le C|v|$.
\item There is a complex Radon measure $\nu$ on $\R\setminus\{0\}$,
with $\sigma$-finite total variation, such that
\begin{equation}\label{eq:increment-class}
 \phi(u)=\int (e^{i\omega u}-1)\,d\nu(\omega),\qquad
 M=\int|\omega|\,d|\nu|(\omega)<\infty.
\end{equation}
The displayed integral is real-valued.
\end{enumerate}
Then Theorem~\ref{thm:main}'s arithmetic-sign conclusion holds for
$\phi$ networks, with constants depending on this fixed activation.
In particular polynomial Frobenius product implies nonuniform
$\PNC^1\subseteq\Lpoly$.
\end{theorem}

The integral in \eqref{eq:increment-class} is absolutely convergent,
since $|e^{i\omega u}-1|\le|\omega u|$. The total variation of $\nu$
itself may be infinite near zero. No efficient sampling algorithm,
computable Fourier coefficients, or quantitative Fourier-tail rate is
assumed. Oddness and analyticity are also unnecessary.

\subsection{The exact scalar transform}

Let $X$ have mean $a$, second moment at most $V>0$, and expected
arithmetic work at most $T\ge1$. At a frequency $\omega$, put
\[
 \lambda=1+\omega^2V,\qquad K\sim\operatorname{Pois}(\lambda),
 \qquad Z=\prod_{j=1}^K(1+i\omega X_j/\lambda),
\]
with independent child calls. The compound-Poisson identity gives
\[
 \E Z=e^{i\omega a},\qquad
 \E|Z-1|^2=e^{\omega^2\E X^2/\lambda}-2\cos(\omega a)+1
       \le C\min(\omega^2V,1).
\]
Expected complex-formula work is $O((1+\omega^2V)T)$.
Write $v(\omega)=\min(\omega^2V,1)$, $c(\omega)=1+\omega^2V$, and
\[
 I=\int\sqrt{vc}\,d|\nu|,\qquad
 J=\int\sqrt{v/c}\,d|\nu|.
\]
Unless $\nu=0$, which is trivial,
\[
 0<J\le I\le\sqrt{2V}M,
\]
because $\sqrt{\min(a^2,1)(1+a^2)}\le\sqrt2a$.
Sample the frequency with density $\sqrt{v/c}/J$ relative to
$|\nu|$, multiply $Z-1$ by the corresponding importance coefficient
and the unit complex phase $d\nu/d|\nu|$, and take its real part.
The resulting unbiased estimator $U$ of $\phi(a)$ has
\begin{equation}\label{eq:increment-potential}
 (\E\,\mathrm{work}(U))\sup_x\E|U(x)|^2
 \le C M^2VT.
\end{equation}
Indeed the two factors are at most $CTI/J$ and $CJI$.
All choices are input-independent. The subtraction of one is part of
the same realized expression; it is not replaced by a separately
sampled term. No higher moment of a child or second Fourier moment
is used. Lemma~\ref{lem:linear} now gives Frobenius propagation with
one fixed activation-dependent factor per retained layer.

\subsection{The short mixed normal form}

The relevant-coordinate and two-ball lemmas use only zero preservation
and one-Lipschitzness, so they hold for $\phi$. There are at most $P^2$
relevant input coordinates, and $j=O(\log(P+2))$ retained prefix
layers give polynomially many separated normalized centers.

Width compression also uses boundedness and local linearization.
Here are the modifications to Appendix~\ref{sec:pf-width-compression}.
At adjacent matrices $W,V$, let $S$ be the preceding norm product,
$P_*$ the current full-product bound, and $e$ the allocated error.
Delete rows with $S\norm{w_i}>B$, where $B=AP_*/e$.
The activation bound $A$ makes their propagated error at most $e$.
For rows with $S\norm{w_i}<\tau$, replace the activation by the
identity, taking
\[
 \tau\le r_0/R,\qquad cP_*\tau^\kappa R^{1+\kappa}\le e.
\]
Their activation-error norm is at most
$c\tau^\kappa R^{1+\kappa}S\norm W_F$, so propagation gives
error at most $e$. Middle-row sampling is unchanged, with ratio
parameter $(B/\tau)^2$, and keeps ordinary $\phi$ neurons.

For the truncated linear branch of nuclear norm $\xi$, replicate
its singular channels $K$ times with the same balanced incoming and
outgoing weights as in that appendix. If $c_0$ is its balancing
scalar, the propagated nonlinear error is at most
\[
 CP_* (c_0\xi S^2/K)^{\kappa/2}R^{1+\kappa}.
\]
To see this, the preactivation norm in one copy is at most
$\sqrt{c_0/K}\sqrt\xi SR$. Apply the local remainder bound,
then the outgoing norm $\sqrt\xi/\sqrt{c_0K}$, sum the $K$
copies, and propagate through the suffix. The existing sampling
estimates bound $c_0\xi S^2$ polynomially in the short-prefix
resources. Choose $K$ polynomially large to obtain error at most
$e$ and keep those preactivations inside $[-r_0,r_0]$.
The exact pair-product calculation is unchanged. Since $\kappa$
is fixed, all retained widths are polynomial in $P$.

Boundedness and the margin give $S_s\le AP\sqrt C$ at retained
prefix layers of width at most $C$. Choose $\epsilon>0$ with inverse
polynomial in $P$ and $\epsilon S_s$ in the local inverse interval.
A constant rail with outputs $\epsilon S_s$ uses first weight
$\phi^{-1}(\epsilon S_1)$ and later weights
\[
 \frac{\phi^{-1}(\epsilon S_s)}{\epsilon S_{s-1}}.
\]
Each is at most a fixed constant times its original layer norm.
The rail therefore increases the prefix product by at most a fixed
constant to the power $j$, hence a polynomial. Normalize the
interface and attach the ordinary-tanh squared-distance classifier
from Appendix~\ref{app:normal}. It uses two hidden layers and
polynomial width, weights and product. No tail estimate or curvature
condition on $\phi$ is needed for this replacement suffix.

For completeness, the physical-weight argument also extends to this
mixed network. Set $\kappa_* =\min(\kappa,2)$. Both activations
have a remainder $C|u|^{1+\kappa_*}$ on a fixed neighborhood of zero.
Replace activations at layers whose original prefix product is less
than $\tau$ by the identity, choosing
\[
 \tau\le r_*/R,\qquad
 CPD\tau^{\kappa_*}R^{1+\kappa_*}\le1/100.
\]
Telescoping controls the output error. Collapse consecutive linear
blocks. Each retained prefix product lies between $\tau$ and
$\max(A,1)P\sqrt C$, so the resulting weight bounds are polynomial.
Double the readout and round using the ordinary one-Lipschitz
perturbation estimate. The resulting mixed network has unit margin,
logarithmic depth, polynomial width and product, polynomial weights,
and $O(\log^2(d+P+2))$ coefficient bits.

\subsection{Finite coefficients without a spectral-tail hypothesis}

An explicit deterministic approximation avoids any hidden regularity
assumption on $\nu$. Let $B\ge2$ be polynomial in $d+P$ and bound
all dimensions, row sums, and preactivation magnitudes in the mixed
normal form, enlarging it by fixed activation-dependent factors.
Its depth is $D=O(\log(d+P+2))$. Choose $\delta$ with
\[
 \log\delta^{-1}=O_\phi(\log^2(d+P+2))
\]
so small that replacing its $\phi$ evaluations by $\delta$-close
functions on $[-2B,2B]$ changes the final output by at most $1/10$.
For example the row-sum estimate
$D\sqrt C[B\max(1,2M)]^D\delta$ suffices.
Boundedness and $\delta\le1$ keep every hybrid preactivation in
this interval inductively.

Define
\begin{align*}
 \phi_{\eta,U}(u)&=a_\eta u+\Rea\int_{\eta\le|\omega|\le U}
                            (e^{i\omega u}-1)\,d\nu(\omega),\\
 a_\eta&=\Rea\int_{0<|\omega|<\eta}i\omega\,d\nu(\omega).
\end{align*}
Taylor's second-order remainder and the bounded exponential increment
imply
\begin{equation}\label{eq:frequency-truncation}
 |\phi(u)-\phi_{\eta,U}(u)|
 \le\tfrac12M\eta|u|^2+2M/U.
\end{equation}
Choose cutoffs making this at most $\delta$ on the interval above.
Both $|\log\eta|$ and $\log U$ are
$O_\phi(\log^2(d+P+2))$. The new activation is at most
$2M$-Lipschitz. This is a deterministic replacement with bounded
bias, not a claim that a conditionally truncated estimator remains
unbiased for $\phi$.

The truncated Fourier branch retains
\eqref{eq:increment-potential}. The linear branch has work $T$ and
moment at most $M^2V$, even if $a_\eta=0$. Combine the two by sampling
branch $i$ with probability proportional to $\sqrt{v_i/t_i}$ for
its positive moment and work bounds. The resulting potential is at
most $C(\sum_i\sqrt{v_it_i})^2\le C_\phi VT$.
Thus logarithmic depth and polynomial Frobenius product still give
polynomial final work--moment potential.

We spell out the scale control. Fix a compact frequency band
$a\le|\omega|\le b$ with positive $|\nu|$ mass; it exists since
$\phi$ is not zero. Choose all truncations to contain it. The
truncated factory then satisfies
\[
 J\ge c_\phi\min(\sqrt V,1/\sqrt V),\qquad
 J\le I\le\sqrt{2V}M.
\]
These estimates, the positive upper moment bound $M^2V$ for the
linear branch, and Lemma~\ref{lem:linear} imply the same logarithmic
recurrence as Lemma~\ref{lem:constants}: at each retained layer
logarithmic magnitudes of deterministic positive quantities and their
inverses increase by at most a fixed factor plus a polynomial.
Nonzero original matrix entries have polynomially bounded logarithmic
reciprocals. The $O(\log(d+P+2))$ layers therefore give polynomial
logarithmic bounds for every internal work, moment and sampling
normalization. No reciprocal of $a_\eta$ is used.

All sampled frequencies are now between $\eta$ and $U$. Their
importance coefficient has modulus
\[
 J\sqrt{\frac{1+\omega^2V}{\min(\omega^2V,1)}},
\]
with polynomial logarithmic size by the preceding bounds. The same
holds for $\omega/\lambda$, the bounded coefficient $a_\eta$,
and all cost-aware neuron coefficients. Average the output estimators
and take $O(d+1)$ groups. The strict-margin, Chernoff, union-bound and
Markov argument of Lemma~\ref{lem:realization} now selects polynomial
total formula size with every constant of modulus at most
$\exp(\poly(d+P+2))$. No frequency-tail event is needed.
Lemma~\ref{lem:round}, denominator clearing, and the single global
real-affine balancing proof apply to these realized expressions.
They give the claimed polynomial-bit integer arithmetic-sign
circuits. This proves Theorem~\ref{thm:fourier-class}.

\subsection{Examples and scope}

A convenient sufficient condition is that $\phi'$ be the Fourier
transform of a finite complex measure $\mu$ with no atom at zero.
Integrating the derivative gives \eqref{eq:increment-class} with
$d\nu=d\mu/(i\omega)$ and $M=\norm\mu_{\mathrm{TV}}$.
The following examples have positive probability derivative measures,
so $M=1$.
\begin{itemize}
\item $\phi(u)=\int_0^u e^{-s^2}\,ds=(\sqrt\pi/2)\operatorname{erf}(u)$:
its derivative is the characteristic function of $N(0,2)$.
It has a cubic local remainder.
\item $\phi(u)=\arctan u$: its derivative $1/(1+u^2)$ is the
characteristic function of the density $e^{-|\omega|}/2$.
It also has a cubic local remainder.
\item $\phi(u)=u/(1+|u|)$: its derivative is continuous and
\[
 (1+|u|)^{-2}=\int_0^\infty t e^{-t}e^{-t|u|}\,dt.
\]
Each $e^{-t|u|}$ is a Cauchy characteristic function, and the mixing
density $te^{-t}$ has mass one. The local remainder is at most
$|u|^2$ and the local inverse is $v/(1-|v|)$. This example is not
analytic at zero.
\end{itemize}
All are bounded, one-Lipschitz, zero-preserving and locally invertible.

The coded activation in the analytic depth counterexample of
\cite{analytic} also
has a finite derivative Fourier measure. A single bump has form
$\delta b_\rho((u-c)/r)$, where
$b_\rho(v)=(1+\rho v)^3e^{-v^2}$ and $|\rho|\le1$.
Its derivative Fourier $L^1$ norm is at most $C\delta/r$, by scaling
and the uniform bound for this fixed-degree Gaussian polynomial.
For its all-derivatives-bounded construction, with
$m_d=M_0+\lceil\log_2(d+2)\rceil$, summing the $2^d$
bumps in band $d$ gives the finite bound
\[
 C\tau\sum_d2^{-(m_d-4)d}<\infty.
\]
The original constant-order construction has the same summability.
Add the Fourier measure of $\operatorname{sech}^2$ and the reflected
perturbation. This is consistent with its learning lower bound:
those teachers have large Frobenius product. Fourier compressibility
of low-product networks is not efficient numerical evaluation of the
activation primitive.

A Bayesian consequence requires a separately established posterior
low-product ceiling for the chosen activation. Theorem~\ref{thm:fourier-class}
does not supply that ceiling or resolve arbitrary conditional controls
at $\gamma=1$. ReLU and GELU are not covered by this bounded-activation
theorem; it makes no claim for all smooth or all Lipschitz activations.
